\section{网络编程}

\begin{question}
	以下关于Java中Socket编程的说法错误的是
	\begin{options}
		\item 用于网络通信
		\item ServerSocket用于服务器端
		\item Socket用于客户端
		\item 不需要处理异常
	\end{options}
	\begin{answer}
		D
	\end{answer}
	\begin{explanation}
		网络编程必须处理 \texttt{IOException} 等异常。
	\end{explanation}
\end{question}

\begin{question}
	以下哪个方法用于发送数据？
	\begin{options}
		\item send()
		\item write()
		\item output()
		\item transmit()
	\end{options}
	\begin{answer}
		B
	\end{answer}
	\begin{explanation}
		使用 \texttt{OutputStream.write()} 发送字节数据。
	\end{explanation}
\end{question}

\begin{question}
	以下关于Java中URL类的说法错误的是
	\begin{options}
		\item 用于表示统一资源定位符
		\item 可以获取网络资源
		\item 不能解析URL
		\item 可以获取URL的各个部分
	\end{options}
	\begin{answer}
		C
	\end{answer}
	\begin{explanation}
		\texttt{URL} 类可以解析 URL 字符串。
	\end{explanation}
\end{question}

\begin{question}
	以下哪个方法用于获取URL的协议部分？
	\begin{options}
		\item getProtocol()
		\item getScheme()
		\item getType()
		\item getProtocolType()
	\end{options}
	\begin{answer}
		B
	\end{answer}
	\begin{explanation}
		\texttt{getScheme()} 返回协议（如 http、https）。
	\end{explanation}
\end{question}

\begin{question}
	以下关于Java中Servlet的说法错误的是
	\begin{options}
		\item 运行在服务器端
		\item 是一个接口
		\item 用于处理HTTP请求
		\item 可以直接被实例化
	\end{options}
	\begin{answer}
		D
	\end{answer}
	\begin{explanation}
		Servlet 由容器管理实例化，开发者不直接 new。
	\end{explanation}
\end{question}

\begin{question}
	以下哪个方法是Servlet的生命周期方法？
	\begin{options}
		\item init()
		\item start()
		\item execute()
		\item process()
	\end{options}
	\begin{answer}
		A
	\end{answer}
	\begin{explanation}
		\texttt{init()}、\texttt{service()}、\texttt{destroy()} 是生命周期方法。
	\end{explanation}
\end{question}

\begin{question}
	以下关于Java中JSP的说法错误的是
	\begin{options}
		\item 本质是一个Servlet
		\item 可以嵌入Java代码
		\item 性能比Servlet高
		\item 用于生成动态网页
	\end{options}
	\begin{answer}
		C
	\end{answer}
	\begin{explanation}
		JSP 首次访问需编译为 Servlet，性能略低于直接使用 Servlet。
	\end{explanation}
\end{question}

\begin{question}
	以下哪个是JSP中的指令？
	\begin{options}
		\item <jsp:include>
		\item <jsp:forward>
		\item <%@ %>
		\item <jsp:param>
	\end{options}
	\begin{answer}
		C
	\end{answer}
	\begin{explanation}
		<%@ page %>、<%@ include %> 是指令（Directive），而 <jsp:include> 是动作（Action）。
	\end{explanation}
\end{question}
